%========================================================================================
%  Seção 15.6 do D2L -- Subword Embedding
%  https://d2l.ai/chapter_natural-language-processing-pretraining/subword-embedding.html
%========================================================================================

%------------------------------------------------
\section{Subword embedding}
%------------------------------------------------

\begin{frame}{Subword Embedding: olhando para dentro das palavras}
    {\footnotesize\color{gray}(cada palavra possui uma estrutura morfológica)}
    \begin{columns}[c]
        \begin{column}{0.49\textwidth}
            \footnotesize
            \begin{itemize}
                \item \textbf{Raízes e sufixos são compartilhados e possuem significado próprio:}
                \begin{itemize}
                    \item \textbf{ajud}-ar $\to$ \textbf{ajud}-ou $\to$ \textbf{ajud}-ando (mesma ação, tempos verbais diferentes)
                \end{itemize}
                \item \textbf{A relação morfológica se repete sistematicamente em pares:}
            \end{itemize}
            \vspace{0.1cm}
            \centering
            \footnotesize
            \begin{tabular}{rcl}
                gat-o $\to$ gat-os & $\approx$ & cã-o $\to$ cã-es {\tiny (plural)} \\
                gat-o $\to$ gat-inho & $\approx$ & cão $\to$ cão-zinho {\tiny (diminutivo)} \\
            \end{tabular}
        \end{column}
        \begin{column}{0.48\textwidth}
            \begin{alertblock}{modelos anteriores (word2vec, GloVe) ignoram isso}
                \scriptsize
                Cada palavra recebe um vetor totalmente \textbf{independente}. Para o modelo, ``gato'' e ``gatinho'' são tão distantes quanto ``gato'' e ``avião''.
            \end{alertblock}
            \begin{alertblock}{consequência prática}
                \scriptsize
                Palavras raras têm vetores de baixa qualidade, e palavras \textbf{fora do vocabulário} (ex.: ``desajudar'') não possuem representação vetorial.
            \end{alertblock}
        \end{column}
    \end{columns}
    \begin{exampleblock}{ideia da seção (ex.: fastText)}
        \scriptsize
        Representar palavras pela \textbf{soma dos vetores dos seus pedaços} (ex.: $\vec{\text{gat}} + \vec{\text{inho}} + \vec{\text{s}}$). Palavras parecidas ou inéditas reaproveitam o conhecimento das partes.
    \end{exampleblock}

    \note{
        Os modelos que vimos até agora, word2vec e GloVe, tratam cada palavra como um bloco indivisível. Mas as palavras têm estrutura morfológica. Raízes e sufixos são compartilhados e têm significado próprio: ``ajudar'', ``ajudou'' e ``ajudando'' são a mesma ação em tempos verbais diferentes. E a relação se repete em pares: de ``gato'' para ``gatos'' é a mesma relação que de ``cão'' para ``cães'' (plural); de ``gato'' para ``gatinho'', a mesma que de ``cão'' para ``cãozinho'' (diminutivo).
        No word2vec e no GloVe, cada palavra tem um vetor totalmente independente. Para o modelo, ``gato'' e ``gatinho'' são tão distantes quanto ``gato'' e ``avião''.
        A consequência prática: palavras raras ficam com vetores de baixa qualidade, e uma palavra fora do vocabulário, como ``desajudar'', não tem representação vetorial.
        A ideia desta seção, por exemplo no fastText, é representar a palavra pela soma dos vetores dos seus pedaços: gat mais inho mais s. Palavras parecidas ou inéditas reaproveitam o conhecimento das partes. Vamos ver duas formas: o fastText e o Byte Pair Encoding.
    }
\end{frame}

\begin{frame}{fastText: a palavra como um saco de n-gramas}
    \small
    \begin{itemize}
        \item fastText \cite{bojanowski2017enriching}: marca o início e o fim com \texttt{<} e \texttt{>} e quebra a palavra em \textbf{n-gramas de caracteres}.
    \end{itemize}
    \vspace{0.1cm}
    \centering
    \begin{tikzpicture}[
        ch/.style={draw=dexpiBlue, fill=dexpiBlueLight, minimum width=0.6cm, minimum height=0.6cm, font=\small\ttfamily, inner sep=0pt},
        ng/.style={draw=dexpiGreenDark, fill=dexpiGreenLight, rounded corners=2pt, font=\footnotesize\ttfamily, minimum height=0.55cm, inner sep=3pt}]
        \foreach \c [count=\k] in {<,w,h,e,r,e,>} {
            \node[ch] (c\k) at (0.6*\k,0) {\c};
        }
        \draw[decorate, decoration={brace, amplitude=4pt, mirror}, thick, dexpiGreenDark] ($(c1.south west)+(0.03,-0.08)$) -- ($(c3.south east)+(-0.03,-0.08)$);
        \draw[decorate, decoration={brace, amplitude=4pt}, thick, dexpiRedDark] ($(c2.north west)+(0.03,0.08)$) -- ($(c4.north east)+(-0.03,0.08)$);
        \node[font=\tiny, text=dexpiGreenDark] at ($(c2.south)+(0,-0.45)$) {janela de 3};
        \node[font=\tiny, text=dexpiRedDark] at ($(c3.north)+(0,0.45)$) {desliza 1 casa};

        \node[ng] (g1) at (6.6,0.55) {<wh};
        \node[ng, right=0.15cm of g1] (g2) {whe};
        \node[ng, right=0.15cm of g2] (g3) {her};
        \node[ng, right=0.15cm of g3] (g4) {ere};
        \node[ng, right=0.15cm of g4] (g5) {re>};
        \node[ng, fill=dexpiLightBlue, draw=dexpiBlue, below=0.35cm of g3] (g6) {<where>};
        \node[font=\tiny, anchor=west] at (g6.east) {\;+ a palavra inteira};
        \draw[-{Stealth}, thick, dexpiBlue] (4.6,0) -- (5.9,0.4);
    \end{tikzpicture}

    \vspace{0.3cm}
    \raggedright
    \begin{itemize}
        \item Os marcadores separam prefixos e sufixos: o n-grama \texttt{her} (dentro de ``where'') é \textbf{diferente} da palavra \texttt{<her>}.
        \item Na prática, usam-se todos os n-gramas com $3 \le n \le 6$.
    \end{itemize}
    \begin{exampleblock}{Por que ajuda}
        \footnotesize
        ``gato'' e ``gatos'' compartilham \texttt{<ga}, \texttt{gat}, \texttt{ato}: os dois vetores passam a ter \textbf{partes em comum}.
    \end{exampleblock}

    \note{
        O fastText, de Bojanowski e colegas, de 2017, é a primeira solução. Ele trata cada palavra como um conjunto de pedaços de caracteres, os n-gramas.
        Primeiro, colocamos um marcador de início, o sinal de menor, e um de fim, o sinal de maior. Depois, deslizamos uma janela de tamanho 3 pela palavra. Para ``where'', saem os n-gramas: menor-w-h, w-h-e, h-e-r, e-r-e e r-e-maior. Também incluímos a palavra inteira, com os marcadores, como um item especial.
        Os marcadores importam: o n-grama ``her'', que aparece no meio de ``where'', é diferente da palavra ``her'' sozinha, que vira menor-h-e-r-maior. Assim o modelo distingue prefixos, sufixos e o meio da palavra.
        Na prática, não usamos só n-gramas de tamanho 3: usamos todos os tamanhos de 3 a 6.
        E por que isso ajuda? Porque ``gato'' e ``gatos'' compartilham vários n-gramas, como menor-g-a, g-a-t e a-t-o. Então os vetores das duas palavras passam a ter partes em comum.
    }
\end{frame}

\begin{frame}{fastText: o vetor da palavra é a soma dos pedaços}
    \small
    \begin{itemize}
        \item Cada subpalavra $g$ tem seu próprio vetor $\mathbf{z}_g$. Com $\mathcal{G}_w$ = subpalavras de $w$:
    \end{itemize}
    \begin{equation*}
        \boxed{\;\vv_w = \sum_{g \in \mathcal{G}_w} \mathbf{z}_g\;}
        \qquad
        \text{ex.: } \vv_{\text{where}} = \mathbf{z}_{\texttt{<wh}} + \mathbf{z}_{\texttt{whe}} + \cdots + \mathbf{z}_{\texttt{<where>}}
    \end{equation*}
    \begin{itemize}
        \item O resto é o \textbf{skip-gram} de sempre: só trocamos $\vv_w$ por essa soma.
        \item O gradiente chega em cada $\mathbf{z}_g$: treinar ``gatos'' também melhora ``gato''.
    \end{itemize}
    \begin{columns}[t]
        \begin{column}{0.48\textwidth}
            \begin{exampleblock}{Ganho}
                \footnotesize
                Palavra \textbf{nova} (``gatinhos''): somamos os vetores dos n-gramas dela, que já foram treinados em outras palavras.
            \end{exampleblock}
        \end{column}
        \begin{column}{0.48\textwidth}
            \begin{alertblock}{Custo}
                \footnotesize
                Há \textbf{muito mais} subpalavras que palavras: mais parâmetros e mais conta por palavra.
            \end{alertblock}
        \end{column}
    \end{columns}

    \note{
        Como o fastText usa esses pedaços? Cada subpalavra $g$ ganha o seu próprio vetor, $\mathbf{z}_g$. O vetor de uma palavra deixa de ser um parâmetro independente: ele passa a ser a soma dos vetores das suas subpalavras. Por exemplo, o vetor de ``where'' é o vetor de menor-w-h, mais o de w-h-e, e assim por diante, mais o da palavra inteira.
        O resto do modelo é o skip-gram que já conhecemos. Só trocamos o vetor da palavra central por essa soma.
        Como o vetor é uma soma, o gradiente chega em cada pedaço. Então, quando o modelo treina com ``gatos'', ele também ajusta os pedaços que ``gatos'' divide com ``gato''. As palavras parecidas aprendem juntas.
        O ganho principal: uma palavra que nunca apareceu no treino, como ``gatinhos'', ainda ganha um vetor, somando os vetores dos seus n-gramas, que já foram treinados em outras palavras.
        O custo: existem muito mais subpalavras do que palavras. O modelo tem mais parâmetros, e calcular o vetor de cada palavra exige somar vários pedaços.
    }
\end{frame}

\begin{frame}{Byte Pair Encoding (BPE)}
    \small
    \begin{itemize}
        \item No fastText, os pedaços têm tamanho \textbf{fixo} ($3$ a $6$) e todos entram no vocabulário.
        \item BPE \cite{sennrich2016neural}: deixa os \textbf{dados} decidirem quais pedaços valem a pena, com tamanhos \textbf{variáveis}.
    \end{itemize}
    \begin{block}{O algoritmo}
        \footnotesize
        \begin{enumerate}
            \item Comece com o vocabulário de \textbf{caracteres} (mais \texttt{\_} para o fim da palavra).
            \item Conte, no corpus, cada \textbf{par de símbolos vizinhos}.
            \item \textbf{Junte} o par mais frequente num símbolo novo.
            \item Repita os passos 2 e 3 um número fixo de vezes.
        \end{enumerate}
    \end{block}
    \begin{exampleblock}{Resultado}
        \footnotesize
        Pedaços frequentes viram símbolos; palavras raras se quebram em pedaços menores. Usado no GPT-2 e no RoBERTa.
    \end{exampleblock}

    \note{
        O fastText usa pedaços de tamanho fixo, de 3 a 6 caracteres, e todos eles entram no vocabulário, o que deixa o vocabulário enorme. O Byte Pair Encoding, ou BPE, de Sennrich e colegas, faz diferente: deixa os dados decidirem quais pedaços valem a pena, e os pedaços podem ter qualquer tamanho.
        O algoritmo é simples. Começamos com um vocabulário de caracteres, mais um símbolo especial, o sublinhado, que marca o fim da palavra. Contamos, no corpus inteiro, quantas vezes cada par de símbolos aparece um ao lado do outro. Juntamos o par mais frequente num símbolo novo. E repetimos isso um número fixo de vezes.
        O resultado: os pedaços que aparecem muito viram símbolos próprios, e as palavras raras se quebram em pedaços menores. É o método usado em modelos de linguagem pré-treinados como o GPT-2 e o RoBERTa.
    }
\end{frame}

\begin{frame}{BPE com números: as primeiras junções}
    \begin{columns}[c]
        \begin{column}{0.42\textwidth}
            \small
            \textbf{Corpus} (palavra: frequência):
            \vspace{0.1cm}

            \centering
            \footnotesize
            \begin{tabular}{lc}
                \toprule
                símbolos & freq. \\
                \midrule
                \texttt{f a s t \_} & 4 \\
                \texttt{f a s t e r \_} & 3 \\
                \texttt{t a l l \_} & 5 \\
                \texttt{t a l l e r \_} & 4 \\
                \bottomrule
            \end{tabular}

            \vspace{0.2cm}
            \raggedright
            Contagem do par (\texttt{t}, \texttt{a}): aparece em ``tall'' e ``taller'':
            $5 + 4 = \textbf{9}$.
        \end{column}
        \begin{column}{0.54\textwidth}
            \centering
            \footnotesize
            \begin{tabular}{clcl}
                \toprule
                junção & par & freq. & novo símbolo \\
                \midrule
                1 & (\texttt{t}, \texttt{a}) & 9 & \nz{\texttt{ta}} \\
                2 & (\texttt{ta}, \texttt{l}) & 9 & \nz{\texttt{tal}} \\
                3 & (\texttt{tal}, \texttt{l}) & 9 & \nz{\texttt{tall}} \\
                4 & (\texttt{f}, \texttt{a}) & 7 & \nz{\texttt{fa}} \\
                5 & (\texttt{fa}, \texttt{s}) & 7 & \nz{\texttt{fas}} \\
                6 & (\texttt{fas}, \texttt{t}) & 7 & \nz{\texttt{fast}} \\
                7 & (\texttt{e}, \texttt{r}) & 7 & \nz{\texttt{er}} \\
                8 & (\texttt{er}, \texttt{\_}) & 7 & \nz{\texttt{er\_}} \\
                9 & (\texttt{tall}, \texttt{\_}) & 5 & \nz{\texttt{tall\_}} \\
                10 & (\texttt{fast}, \texttt{\_}) & 4 & \nz{\texttt{fast\_}} \\
                \bottomrule
            \end{tabular}
        \end{column}
    \end{columns}
    \vspace{0.3cm}
    \begin{exampleblock}{Depois de 10 junções}
        \small
        \texttt{fast\_} \quad \texttt{fast er\_} \quad \texttt{tall\_} \quad \texttt{tall er\_}
        \hfill {\footnotesize o sufixo \texttt{er\_} virou um símbolo próprio}
    \end{exampleblock}

    \note{
        Vamos rodar o BPE num exemplo do livro. O corpus tem quatro palavras: ``fast'' aparece 4 vezes, ``faster'' 3, ``tall'' 5 e ``taller'' 4. Cada palavra começa quebrada em caracteres, com o sublinhado no fim.
        Contamos os pares vizinhos, levando em conta a frequência das palavras. O par t-a aparece em ``tall'' e em ``taller'', então a contagem é 5 mais 4, 9. É o par mais frequente, e vira o símbolo ``ta''.
        Recontamos e repetimos. Na tabela, as dez junções em ordem: primeiro se forma ``tall'', com frequência 9; depois ``fast'', com 7; depois o sufixo ``er'' e ``er'' com o fim de palavra, também com 7; e por fim ``tall'' e ``fast'' com o fim de palavra.
        Depois das dez junções, as palavras ficam assim: ``fast'' vira um símbolo só, ``faster'' vira ``fast'' mais ``er'', ``tall'' vira um símbolo só, e ``taller'' vira ``tall'' mais ``er''. Repare que o sufixo ``er'' virou um símbolo próprio: o algoritmo descobriu sozinho um pedaço com significado.
    }
\end{frame}

\begin{frame}{BPE: segmentando palavras que nunca vimos}
    \small
    \begin{itemize}
        \item Para uma palavra nova, pegamos sempre o \textbf{maior pedaço conhecido} a partir do início (busca gulosa).
    \end{itemize}
    \vspace{0.2cm}
    \centering
    \begin{tikzpicture}[
        tok/.style={draw=dexpiBlue, fill=dexpiBlueLight, rounded corners=2pt, font=\small\ttfamily, minimum height=0.6cm, inner sep=4pt},
        kn/.style={tok, draw=dexpiGreenDark, fill=dexpiGreenLight},
        lab/.style={font=\small, anchor=east}]
        \node[lab] at (-0.3,0) {tallest\_ $\;\Rightarrow$};
        \node[kn, anchor=west] (a1) at (0,0) {tall};
        \node[tok, right=0.1cm of a1] (a2) {e};
        \node[tok, right=0.1cm of a2] (a3) {s};
        \node[tok, right=0.1cm of a3] (a4) {t};
        \node[tok, right=0.1cm of a4] (a5) {\_};

        \node[lab] at (-0.3,-1) {fatter\_ $\;\Rightarrow$};
        \node[kn, anchor=west] (b1) at (0,-1) {fa};
        \node[tok, right=0.1cm of b1] (b2) {t};
        \node[tok, right=0.1cm of b2] (b3) {t};
        \node[kn, right=0.1cm of b3] (b4) {er\_};
    \end{tikzpicture}

    \vspace{0.3cm}
    \raggedright
    \begin{itemize}
        \item Em verde, símbolos aprendidos nas junções; em azul, caracteres soltos.
        \item ``tallest'' reaproveita \texttt{tall}; ``fatter'' reaproveita o sufixo \texttt{er\_}.
        \item Nenhuma palavra fica sem representação: no pior caso, ela se quebra em \textbf{caracteres}.
    \end{itemize}
    \begin{exampleblock}{Na prática}
        \footnotesize
        Os vetores (ou o modelo de linguagem) são treinados sobre esses símbolos, e não sobre palavras inteiras.
    \end{exampleblock}

    \note{
        Com o vocabulário de símbolos aprendido, como quebramos uma palavra que nunca vimos? Com uma busca gulosa: a partir do início da palavra, pegamos sempre o maior pedaço que está no vocabulário, e continuamos do ponto onde ele termina.
        ``tallest'' não estava no corpus. O maior pedaço conhecido no começo é ``tall''. O que sobra, e-s-t e o fim de palavra, não forma nenhum símbolo aprendido, então fica em caracteres.
        ``fatter'' também é nova. Ela começa com ``fa'', porque ``fat'' não é um símbolo; depois vêm dois t soltos; e termina com o sufixo ``er'' com fim de palavra, que o algoritmo aprendeu.
        Repare que as duas palavras novas reaproveitam pedaços aprendidos de outras palavras. E nenhuma palavra fica sem representação: no pior caso, ela se quebra até os caracteres.
        Na prática, os vetores ou o modelo de linguagem são treinados sobre esses símbolos, e não sobre as palavras inteiras.
    }
\end{frame}

